% Altimeter Pipeline — Validation Results Section
% Drop into your paper's results section.

\subsection{Pipeline validation}

\subsubsection{Cross-site QC performance}

The pipeline was validated across five representative deployments spanning
all three sites, both instrument types (AA400 altimeter and EA400
echosounder), and three sampling modes (continuous, burst-300, burst-60).
Table~\ref{tab:alt_validation} summarizes the QC performance.

\begin{table}[h]
\centering
\caption{Pipeline validation across representative deployments.
IQR = median within-burst interquartile range.}
\label{tab:alt_validation}
\begin{tabular}{lcccccc}
\toprule
Deployment & Instrument & Bursts & Rejected & BL range & IQR & Valid\% \\
\midrule
SIO24 (6m, continuous) & AA400 & 257 & 34 & 71~mm & 4.6~mm & 97\% \\
SIO24 (6m, burst-300) & AA400 & 1531 & 6 & 57~mm & 5.1~mm & 97\% \\
TP24 (10m, echosounder) & EA400 & 1719 & 70 & 94~mm & 7.2~mm & 94\% \\
TP25 (10m, burst-60) & AA400 & 6883 & 275 & 128~mm & 3.8~mm & 97\% \\
SOL25 (7m, tilted) & EA400 & 2600 & 490 & 99~mm & 7.3~mm & 88\% \\
\bottomrule
\end{tabular}
\end{table}

The AA400 altimeters consistently exhibited lower within-burst IQR
(\SIrange{3.8}{5.1}{\milli\meter}) compared to the EA400 echosounders
(\SIrange{7.2}{7.3}{\milli\meter}), reflecting the simpler measurement
geometry and more focused acoustic beam of the dedicated altimeter.
The percent of valid samples within accepted bursts exceeded 88\%
across all deployments. The Solana Beach deployment demonstrated that
the tilt-deviation mask correctly preserves data from a tilted
instrument (baseline tilt \SI{13.2}{\degree}) while rejecting periods
of transient disturbance.

\subsubsection{Firmware error recovery}

The Torrey Pines 5m altimeter (Nov 2023 -- Feb 2024) was affected by
a firmware bug that overwrote valid altitude data with date-string
error messages after memory filled ($\sim$113~days). The two-level
despiking approach recovered the underlying morphological signal:
the per-ping phase-space despike removed individual corrupted samples
within bursts, and the iterative burst-level despike (5~passes)
removed 1472~additional corrupted bursts. The resulting time series
clearly shows the \SI{700}{\milli\meter} accretion event caused by
the December 28, 2023 storm, consistent with the subaerial beach
erosion observed in concurrent surveys.

\subsubsection{Independent validation against beach surveys}

Echosounder bed level measurements were independently validated
against beach morphology surveys conducted by the Scripps Coastal
Processes Group using jetski-mounted RTK GPS systems. Survey profiles
were interpolated to the instrument cross-shore position at each
survey date.

For single-deployment validation at Solana Beach (March--June 2025),
the RMSE of the bed level change between consecutive survey pairs was
\SI{22}{\milli\meter} with a bias of \SI{-4}{\milli\meter}. This is
comparable to the spatial variability inherent in the survey comparison
itself (survey-to-survey $\sigma = \SI{21}{\milli\meter}$), indicating
that the instrument accuracy is at least as good as can be resolved by
the available ground truth.

For multi-deployment validation spanning the full record at each site,
the cross-deployment baseline anchoring introduced additional
uncertainty. Torrey Pines 10m (survey-anchored, 3~pairs) showed
RMSE = \SI{61}{\milli\meter}. The Solana Beach multi-deployment record
showed larger RMSE (\SI{92}{\milli\meter}), attributable to physical
relocation of the echosounder pipe during maintenance (re-jetting),
which changed the spatial point of measurement between deployments.

\subsubsection{L4 PUV merge prototype}

The L4 merged product was prototyped on SIO24A (SIO Pier, April 2024),
matching 131~altimeter bursts to PUV L2 wave spectral segments within
$\pm$\SI{5}{\minute} (median offset \SI{2.5}{\minute}). At \SI{6}{\meter}
depth, all matched bursts exceeded the sediment mobilization threshold
($H_s$ = \SIrange{0.55}{1.59}{\meter}, $\tau_b$ = \SIrange{0.20}{1.84}{\pascal}).
Initial scatter plots of $dz/dt$ versus $\tau_b$ show symmetric scatter
around zero, indicating roughly equal accretion and erosion rates at all
stress levels. Threshold identification requires data from deeper sites
(Torrey Pines 10~m, 15~m) where not all conditions exceed the mobilization
threshold.
